《从零构建推理模型》是一本动手实践指南，通过亲手实现核心技法，帮助读者理解现代推理型 LLM 的工作原理。本书并不从零开始训练一个基础模型，而是从一个预训练的开源基础 LLM 出发，一步步为它扩展推理能力。在此过程中，你将学习如何用基础模型生成文本、用验证器评估答案、通过推理时扩展提升性能、用强化学习进行训练，以及把推理行为蒸馏到更小的模型中。读完本书，你将理解各种主要推理方法如何在一个实际的 LLM 开发工作流中相互配合，以及如何构建属于自己的、小型但可用的推理系统。

本书采用代码优先的方式。“从零”指的是由你自己实现这些推理方法，从而让底层机制变得清晰透明。如果你对基础 LLM 如何构建与预训练的细节感兴趣，我此前出版的《从零构建大语言模型》（\href{http://mng.bz/M96o}{http://mng.bz/M96o}）提供了这方面的基础，不过阅读本书并不需要先读那本书。

\section*{本书的读者对象}
\addcontentsline{toc}{section}{本书的读者对象}

《从零构建推理模型》适合机器学习爱好者、工程师、研究者、学生和软件开发人员阅读，他们希望切实理解推理模型的工作原理，以及如何实现其背后的主要技法。本书写给那些不满足于高层描述、想看到这些系统如何在代码中被构建、评估和改进的读者。

最重要的前提是扎实的 Python 编程经验。有机器学习、深度学习或大型语言模型方面的基础会有所帮助，但不必是专家。本书力求让有动力的初学者也能上手，同时对有经验的从业者来说也具备足够的技术深度。

本书会频繁使用数学示例，因为这类示例在推理研究中既实用，又便于自动验证。不过并不需要高等数学。熟悉基础代数、对向量和矩阵有基本的适应能力，会有助于理解部分代码实现。

有一些 PyTorch 经验或许能让你更快地读懂代码，但这并不是必需的。本书的重点是理解推理方法并清晰地实现它们，而不是掌握深度学习框架的每一个细节。读过《从零构建大语言模型》的读者可能会认出一些更宏观的 LLM 概念，但本书内容完全自成一体。

\section*{本书的组织结构：一张路线图}
\addcontentsline{toc}{section}{本书的组织结构：一张路线图}

本书围绕一条实用的推理模型开发流水线来组织。我们先从一个常规的预训练 LLM 出发，加入评估以便衡量进展，然后探索两条改进推理的主要途径：推理时方法和额外的训练。由于后面的章节会直接建立在前文代码与思想之上，本书最好按顺序阅读。

第 1 章介绍推理在 LLM 语境中的实际含义。它解释推理模型与常规 LLM 有何不同，从较高的层面回顾标准的 LLM 训练流程，并介绍用于改进推理的主要方法。

第 2 章确立起点：加载并使用一个预训练的基础 LLM。内容涵盖编程环境、分词、逐步文本生成，以及键值缓存和模型代码编译等实用加速手段。

第 3 章加入评估。它展示如何可靠地提取最终答案、对照参考解答验证正确性，并为面向数学的推理任务构建一条评估流水线。

有了这些基础，第 4 章和第 5 章聚焦推理时扩展，即在不改变模型权重的情况下改进推理。第 4 章涵盖基于提示的推理、多样化解码和自洽性。第 5 章沿着这个方向继续展开，介绍自我修正、简单的评分策略、置信度估计以及迭代式答案改进。

第 6 章到第 8 章实现训练时方法。第 6 章介绍面向推理模型的强化学习，重点是可验证奖励的强化学习（RLVR）和组相对策略优化（GRPO）。第 7 章在该实现的基础上，介绍实用的监控手段、常见的训练失败模式、奖励劫持，以及 KL 正则化和回复格式奖励等有用的扩展。第 8 章以蒸馏结束正文，展示如何创建教师模型生成的推理数据集、训练一个更小的学生模型，并评估蒸馏结果。

附录是对正文各章的补充。附录 A 提供参考文献与延伸阅读，附录 B 汇总练习解答，附录 C 给出全书使用的 Qwen3 源代码，附录 D 讨论更大的 LLM 变体，附录 E 涵盖批处理与面向吞吐量的执行，附录 F 综述更多 LLM 评估方法，附录 G 展示如何围绕模型构建一个聊天界面。

\section*{关于代码}
\addcontentsline{toc}{section}{关于代码}

为便于跟随练习，本书的代码示例可从 Manning 图书页面 \href{https://www.manning.com/books/build-a-reasoning-model-from-scratch}{https://www.manning.com/books/build-a-reasoning-model-from-scratch} 以及配套的 GitHub 仓库 \href{https://github.com/rasbt/reasoning-from-scratch}{https://github.com/rasbt/reasoning-from-scratch} 获取。练习的解答汇总在附录 B 中，相应的代码文件可从仓库的各章文件夹中链接到。本书的 liveBook（在线）版本提供了可直接运行的代码片段，地址为 \href{https://livebook.manning.com/book/build-a-reasoning-model-from-scratch}{https://livebook.manning.com/book/build-a-reasoning-model-from-scratch}。

本书包含许多源代码示例，既有编号的代码清单，也有与正文混排的代码。这两种情况下，源代码都采用 \texttt{fixed-width} \texttt{font like} \texttt{this} 这样的等宽字体排版，以与普通正文区分。

在许多情况下，原始源代码经过了重新排版；我们添加了换行并调整了缩进，以适应书中可用的页面空间。此外，当代码在正文中有所说明时，代码清单里的注释往往已被删除。许多代码清单配有代码注释，用以突出重要概念。

本书的一个关键目标是易用性，因此代码示例经过精心设计，可以在普通笔记本电脑上高效运行，无需任何特殊硬件。但如果你有可用的 GPU，某些章节会提供有用的建议，说明如何扩大数据集和模型规模，以利用这份额外的算力。

正文各章的代码力求易读，并能在消费级硬件上运行。推理与评估章节在 CPU 和 GPU 上都能很好地运行，而面向训练的章节如果希望复现更大规模的实验，则会受益于 GPU。全书都使用 PyTorch，并搭配预训练的开源 Qwen3 模型；附录 C 给出了底层的 Qwen3 实现，供想更细致地查看该代码的读者参考。

\section*{liveBook 讨论论坛}
\addcontentsline{toc}{section}{liveBook 讨论论坛}

购买《从零构建推理模型》即可免费使用 liveBook——Manning 的在线阅读平台。借助 liveBook 独有的讨论功能，你可以针对全书、或针对特定的节或段落添加评论。做笔记、提出并回答技术问题、从作者和其他用户那里获得帮助，都非常方便。要访问论坛，请前往 \href{https://livebook.manning.com/book/build-a-reasoning-model-from-scratch/discussion}{https://livebook.manning.com/book/build-a-reasoning-model-from-scratch/discussion}。

Manning 对读者的承诺是提供一个场所，让读者之间、以及读者与作者之间能够进行有意义的对话。这并不构成作者必须参与一定次数讨论的承诺，作者对论坛的贡献始终是自愿的（且没有报酬）。建议你不妨向作者提一些有挑战性的问题，以免他们失去兴趣！只要本书仍在印行，论坛以及过往讨论的存档都可以从出版社网站上访问。

\section*{其他在线资源}
\addcontentsline{toc}{section}{其他在线资源}

想了解最新的人工智能与大型语言模型（LLM）研究趋势吗？

\begin{itemize}
\item 欢迎访问我的博客 \href{https://magazine.sebastianraschka.com}{https://magazine.sebastianraschka.com}，我会在那里定期讨论最新的人工智能研究，重点关注 LLM。
\end{itemize}

想快速掌握深度学习与 PyTorch？

\begin{itemize}
\item 我在自己的网站 \href{https://sebastianraschka.com/teaching}{https://sebastianraschka.com/teaching} 上提供了若干免费课程。这些资源可以帮助你快速掌握最新技术。
\end{itemize}

想找与本书相关的附加材料？

\begin{itemize}
\item 访问本书的 GitHub 仓库 \href{https://github.com/rasbt/reasoning-from-scratch}{https://github.com/rasbt/reasoning-from-scratch}，可以找到供你补充学习的额外资源与示例。
\end{itemize}
